\subsection{SCALES OF COMPARISON}

Let E be a set filtered by a filter with base F, and K a non-discrete valued field (most often K = R or K = C). In the set of functions in \(\mathcal{H}(\mathfrak{F}, \mathrm{K})\) not equivalent to 0 modulo \(R_{\infty}\) (that is, those for which in every set in F there is at least one point where the function does not vanish), the relation `` \(f \ll g\) or f = g'' is an order relation.

DEFINITION 1. One says that a subset E of \(\mathcal{H}(\mathfrak{F},\mathrm{K})\) formed of functions not equivalent to 0 modulo \(R_{\infty}\) is a comparison scale when E is totally ordered by the relation `` \(f \ll g\) or f = g''.

In other words, if \(f\) and \(g\) are functions in \(\mathcal{E}\) then one (and only one) of the relations \(f \ll g\) , \(g \ll f\) , \(f = g\) always holds. It follows that on \(\mathcal{E}\) the relation \(f \asymp g\) (and a fortiori \(|f| \sim a |g|\) , where \(a\) is a number \(> 0\) ) implies \(f = g\) .

Every subset of a comparison scale is clearly also a comparison scale.

Examples. 1) For x real and tending to \(+\infty\) the set of functions \(x^{\alpha}\) ( \(\alpha\) an arbitrary real number) is a comparison scale. The same is true for the functions \((x - a)^{\alpha}\) when F is the set of open intervals with left-hand endpoint a.

\begin{enumerate}
\def\labelenumi{\arabic{enumi})}
\setcounter{enumi}{1}
\item
  For \(z\) complex tending to \(\infty\) the set of functions \(z^n\) ( \(n\) a rational integer) is a comparison scale; so are the functions \((z - a)^n\) when \(\mathfrak{F}\) is the trace on the complement of the point \(a \in \mathbf{C}\) of the filter of neighbourhoods of this point.
\item
  Let F be normed space; the family of functions \(\|\mathbf{x}-\mathbf{a}\|^{\alpha}\) ( \(\alpha\) an arbitrary real number) is a comparison scale when F is the trace on the complement of a of the filter of neighbourhoods of this point. Note that if p and q are two distinct norms on F, the union of the two comparison scales \(\left(p(\mathbf{x}-\mathbf{a})\right)^{\alpha}\) and \(\left(q(\mathbf{x}-\mathbf{a})\right)^{\alpha}\) is not in general a comparison scale.
\item
  For \(x\) real tending to \(+\infty\) the family \(\mathcal{E}\) of functions of the form \(\exp(p(x))\) , where \(p\) runs through the set of polynomials without constant term (with real coefficients), is a comparison scale: it suffices to remark that the quotient of two functions in \(\mathcal{E}\) is again in \(\mathcal{E}\) , and that a function \(\exp(p(x))\) must tend either to 0 or \(+\infty\) if \(p \neq 0\) ; indeed, \(p(x) \sim \alpha x^n\) where \(n > 0\) and \(\alpha \neq 0\) ; if \(\alpha > 0\) then \(p(x) > \frac{1}{2}\alpha x^n\) for \(x\) sufficiently large; if \(\alpha < 0\) then \(p(x) < \frac{1}{2}\alpha x^n\) for \(x\) sufficiently large; in the first case \(\exp(p(x))\) tends to \(+\infty\) , in the second to 0.
\item
  For x real tending to \(+\infty\) the set E of functions of the form \(x^{\alpha}(\log x)^{\beta}\) (defined for x \textgreater{} 1), where \(\alpha\) and \(\beta\) are arbitrary real numbers, is a comparison scale. Indeed, here again the quotient of two functions in E is a function in E; it is enough to show that if \(\alpha\) and \(\beta\) are not both zero then \(x^{\alpha}(\log x)^{\beta}\) tends to 0 or \(+\infty\) ; this is clear if \(\alpha = 0, \beta \neq 0\) ; if \(\alpha > 0\) one has \((\log x)^{-\beta} \ll x^{\alpha}\) , and if \(\alpha < 0\) one has \((\log x)^{\beta} \ll x^{-\alpha}\) for any \(\beta\) , whence the proposition.
\end{enumerate}

Note that this last comparison scale is a totally ordered set (for the relation `` \(f \ll g\) or \(f = g\) '') whose order structure is isomorphic to the lexicographic order on \(R^{2}\) (Set Theory, III, p.~157); recall that in this structure the relation \((\alpha, \beta) < (\gamma, \delta)\) means `` \(\alpha < \gamma\) , or \(\alpha = \gamma\) and \(\beta < \delta\) '').

Similarly, the scale formed by the functions \(\exp(p(x))\) , where p runs through the set \(P_{0}\) of polynomials with no constant term, has order structure isomorphic to the order structure of \(P_{0}\) , in which the relation p \textless{} q implies that the dominant term of the polynomial q - p has a coefficient \textgreater{} 0 (cf.~Alg., VI. 19, Example 2).

Let \(\varphi\) be a map from a set \(F\) into \(E\) , such that \(\overset{-1}{\varphi}(\mathfrak{F})\) is a filter base on \(F\) . If \(\mathcal{E}\) is a comparison scale on \(E\) (for the filter base \(\mathfrak{F}\) ) then the functions \(f \circ \varphi\) , as \(f\) runs through \(\mathcal{E}\) , form a comparison scale on \(F\) (for the filter base \(\overset{-1}{\varphi}(\mathfrak{F})\) ).

